%%%PAGE 270 rot=0 legibility=good summary=零散方法笔记页：RC充放电、平抛对角线、以M为参考系、图象法面积、功率公式与E-t图斜率、极限法、特殊值法、斜面摩擦力做功及v-t图比较
%%%BLOCK id=270-01 ch=09 sec=电容器·充放电 kind=note alt=10 cont=none y=0.05-0.09
$R$越小\quad 电容充放电较快
%%%END
%%%BLOCK id=270-02 ch=04 sec=平抛运动·空间极值 kind=note alt=none cont=none y=0.09-0.12
平抛在空间中对角线处最远
%%%END
%%%BLOCK id=270-03 ch=05 sec=卫星·追及相遇 kind=method alt=04 cont=none y=0.125-0.20
%FIG p270_a 0.12 0.125 0.29 0.17
\notefig[0.25]{figures/p270_a.png}
以$M$为参\ $N$做圆周运动 % 原稿“参”字后未写完（疑为“参考系”）
%%%END
%%%BLOCK id=270-04 ch=06 sec=动能定理·图像法 kind=method alt=90 cont=none y=0.22-0.37
\textbf{图象法}——面积
\begin{itemize}
\item $F$--$x$
\item $a$--$x$
\item $\mu$--$x$
\item $p$--$\unclear{V}$
\end{itemize}
动能Th\ 多过程曲线运动
%%%END
%%%BLOCK id=270-05 ch=06 sec=功率·公式 kind=formula alt=none cont=none y=0.39-0.54
$P$ 定义式\quad $P=\dfrac{W}{t}$
\[
P=F\cdot v\cos\alpha
\]
\[
\bar P=F\bar v
\]
\[
P_{\text{瞬}}=FV_{\text{瞬}}
\]
%%%END
%%%BLOCK id=270-06 ch=06 sec=功率·E-t图像斜率 kind=knowledge alt=none cont=none y=0.38-0.57
\textbf{$E$--$t$图斜率为对应功率}
\begin{itemize}
\item $E_k$--$t$\quad $P_{\text{合外力}}$
\item $E_{P_G}$--$t$\quad $P_G$
\item $E$--$t$\quad $P_{\text{其他}}$ 如 $P_f$ 等
\end{itemize}
%FIG p270_b 0.375 0.415 0.60 0.568
\notefig[0.25]{figures/p270_b.png}
图旁文字：受力水平；\ $P_{G\max}$
%%%END
%%%BLOCK id=270-07 ch=07 sec=动量与能量综合·极限法 kind=method alt=90 cont=none y=0.57-0.61
\textbf{极限法}
\[
V=\sqrt{\frac{2MgR}{M+R}}\qquad M\gg m\to V=\sqrt{2gR}
\] % CHECK: 分母末项原稿像“R”，物理上应为m
%%%END
%%%BLOCK id=270-08 ch=90 sec=解题方法·特殊值法 kind=method alt=none cont=none y=0.61-0.64
\textbf{特殊值法}\quad 有$k$、\unclear{$n$}、类参数、比值、直接赋值。
%%%END
%%%BLOCK id=270-09 ch=06 sec=功能关系·摩擦力做功与v-t图像 kind=method alt=none cont=none y=0.64-0.98
%FIG p270_c 0.055 0.64 0.27 0.75
\notefig[0.25]{figures/p270_c.png}
\begin{align*}
N&=mg\cos\alpha+F\sin\alpha\\
f&=\mu mg\cos\alpha+\mu F\sin\alpha\\
W_f&=-fS=-\mu mgS\cos\alpha-\mu F\sin\alpha=-\mu mgx-\mu Fy % 原稿此处等号写成“÷”形；CHECK: 第二项按x、y的换算应为μFS sinα，原稿似无S
\end{align*}
$\circ$\ $E_{k1}=E_{k2}$ % 原稿“=”与“Ek₂”之间有一小笔画，疑为连笔

\quad $t_1>t_2$

%FIG p270_d 0.30 0.84 0.53 0.98
\notefig[0.25]{figures/p270_d.png}
%%%END
%%%PAGEEND 270
